\section{Related Work}
\label{sec:related}

\paragraph{Proximal Policy Optimization.}
PPO~\citep{schulman2017ppo} is the workhorse on-policy algorithm for deep RL,
using clipped importance-sampling ratios to constrain policy updates.  Earlier
trust-region methods include TRPO~\citep{schulman2015trpo}.  Our work retains
the PPO framework but changes the granularity at which ratios and clipping are
applied.

\paragraph{Generative Models as Policies.}
Diffusion models~\citep{ho2020ddpm,song2021scorebased} have been adopted as
expressive policy classes.  \citet{chi2023diffusion} introduced Diffusion
Policy for imitation learning in robotics.  \citet{black2024ddpo} proposed DDPO,
which fine-tunes diffusion models with policy-gradient methods by treating each
denoising step as an RL action.  Flow-matching
policies~\citep{lipman2023flow,liu2023flow} replace the diffusion SDE with a
simpler ODE-based transport, offering faster sampling and simpler training
objectives.

\paragraph{Stochastic Interpolants and Stochastic Flow Matching.}
\citet{albergo2023stochastic} introduced stochastic interpolants, providing a
unified framework that generalizes both diffusion and flow matching by
interpolating between source and target distributions with optional stochasticity.
\citet{albergo2023building} extended this to build normalizing flows with
stochastic dynamics.  Our stochastic flow policy draws on this framework,
retaining per-step stochasticity so that each conditional
$p_\theta(\mathbf{z}_{k-1}\mid\mathbf{z}_k)$ has tractable Gaussian density.

\paragraph{RL Fine-Tuning of Generative Policies.}
Several works have applied RL to fine-tune generative policies.
\citet{black2024ddpo} compute per-step importance weights within DDPO.
\citet{fan2024dpok} use KL-regularized RL objectives for diffusion fine-tuning.
\citet{clark2024directly} study direct reward fine-tuning of diffusion models.
Most recently, DPPO~\citep{ren2025dppo} formulates diffusion policy optimization
as a two-level MDP and applies PPO with per-step structure, while
ReinFlow~\citep{hu2025reinflow} fine-tunes flow-matching policies with PPO using
injected stochasticity.  Our work is complementary: rather than proposing a new
algorithm, we conduct a controlled comparison of \emph{how} PPO credit and
clipping should be distributed across internal steps, testing uniform, learned
global, and state-dependent weighting schemes that these prior works do not
systematically evaluate.

\paragraph{Credit Assignment in RL.}
Temporal credit assignment is a classical challenge in
RL~\citep{sutton2018rl,minsky1961steps}.  Hindsight credit assignment
methods~\citep{harutyunyan2019hindsight,hung2019optimizing} attempt to
redistribute reward to relevant earlier decisions.  Our per-step weighting can
be viewed as a form of \emph{intra-action} credit assignment: distributing the
task-level advantage across the internal generative steps that produced the
action.

\paragraph{Positioning.}
Our work bridges generative-model fine-tuning and PPO implementation design.
DPPO~\citep{ren2025dppo} and ReinFlow~\citep{hu2025reinflow} demonstrate that
per-step PPO is viable for generative policies, but neither systematically
compares credit-assignment weighting schemes.  Our contribution is this
controlled comparison: we hold the policy architecture and PPO framework fixed
and vary only whether clipping and credit are applied holistically, per-step
uniformly, or per-step with learned (global or state-dependent) weights.
